\paragraph{Two translations of the same picture.}
We tested whether a capable text translator could teach a small image writer by example. For 511 training passages the student saw both the archived modern gloss and an alternative gloss supplied by Qwen3-VL. A matched control saw each archived gloss twice. The two students wrote freely from 1,024 pictured passages in 519 unseen chapters. Relative to archive replay, supplementary supervision increased the total errors by 587 glyph edits against the archive and by 94 against the teacher. The corresponding paired 95\% intervals for the change in error rate were $[-1.03,5.91]$ and $[-2.82,3.59]$ percentage points. Changing the pictured source to another passage of identical length increased the supplemented student's archive edit count from 22,338 to 27,961. Thus the source influenced the output, even though neither student supplied a faithful modern translation in the eighteen-passage reading. The phrase composer produced three, the pretrained text-fed teacher fourteen and the archive fifteen (Fig.~\ref{fig:paired-student}).

\paragraph{Choosing a glyph shared by several routes.}
The decoder could copy a visible character at multiple positions or reproduce precisely the same ink from its exemplar memory. Its original decision selected a route before comparing the visible alternatives. With the weights fixed, pooling probability assigned to identical ink raised first-glyph matches from 393 to 451 of 1,024 against the archive and from 530 to 629 of 1,023 against the teacher. Against the archive, 69 incorrect first marks became correct and eleven correct marks were lost. Applied throughout the answer, pooling reduced archive edits from 22,338 to 19,793 and teacher edits from 21,750 to 19,242. The new decoder ended 1,015 answers within the length limit and wrote 26,015 cells, compared with 1,010 ends and 27,774 cells under the original decoder. Both produced two exact archive translations requiring a change. A second masked reading found zero faithful pooled answers on the same eighteen passages. These results locate a physical decoding error. Fixing it preserves more of the desired ink but leaves sentence meaning unresolved.
